\scanpage{81}
\mspage{8}
a point $x\in X_K$ is entirely determined by knowledge of \deleted{$f_{x_0,x}$}
\[
f_{x_0,x}\in H^1_{\mathrm{dr}}(\Gamma,\widehat\Pi_1\underbrace{(X,x_0)}_{\widehat\Pi(x_0)}),
\]
and a fortiori by knowledge of the \unsure{\inserted{$1$-cocycle}}
\[
f_l:\Gamma\longrightarrow\widehat\Pi^{\unsure{\text{ét}}}(x_0)\qquad\text{for a single}
\]
$l:x_0\longrightarrow x$ --- or what amounts to the same, by knowledge of $g_{l^{-1}\deleted{\unsure{x_0}}}\in H^1_g(\Gamma,\widehat\Pi_1(x_0))$, \inserted{where $l^{-1}:x\longrightarrow x_0$.}
\margin{$f_l=(g_{l^{-1}})^{-1}$}

\unsure{4|\P} The ``explicit'' determination of an $f_l$ --- \deleted{of a $g_{l^{-1}}$} (which thus characterizes $x$ in the case considered) is a very interesting problem, and surely deep. As \deleted{\unsure{one} \unsure{a}} \inserted{the} group $\Gamma$ of course is far from being known and at hand, one must moreover specify what one may well be able to call a ``determination'' (or even explicit determination!) of $f_l$. For this, one must surely consider that a $u\in\Gamma$ is specified by the invariants $f_{l_j}(u)\in\widehat\Pi_1(x_j)$ or $g_{l_j}(u)\in\widehat\Pi_1(y_j)$, associated above with a \inserted{finite} system of generators \inserted{$l_j:x_j\longrightarrow y_j$} of the groupoid $\Pi_1(X,I)$ ($j\in J$, set of indices for the generating arrows, $x_j,y_j\in I$),
\margin{80}
